% Chapter 11: Quasi-Monte Carlo.
\providecommand{\Corr}{\operatorname{Corr}}

Monte Carlo integrates with random points, and its error $\sigma/\sqrt n$ (\cref{thm:mc-clt})
is the price of that randomness: independent points leave gaps and clusters by chance.
\emph{Quasi-Monte Carlo} (QMC) replaces them by deterministic points built to fill the unit
cube as evenly as possible. For integrands that are smooth enough, the error then decreases
almost like $1/n$ instead of $1/\sqrt n$, in any fixed dimension. The catch lies in ``smooth
enough'' and ``fixed dimension'': option payoffs have kinks and jumps, the map from
uniforms to normals is unbounded, and a path has as many dimensions as it has dates. This
chapter separates what is proved from what is observed, and measures the rest.

\Cref{sec:qmc-discrepancy} measures evenness by the discrepancy and bounds the error by the
Koksma--Hlawka inequality. \Cref{sec:qmc-sequences} builds the classical low-discrepancy
sequences of van der Corput, Halton and Sobol'. \Cref{sec:qmc-rqmc} randomises them, which
restores unbiasedness and the error estimates of \cref{ch:06-monte-carlo}.
\Cref{sec:qmc-effective} explains, through the effective dimension, why the Brownian-bridge
and principal-component constructions of \cref{ch:07-path-simulation} make QMC so effective
for Asian options.

Throughout, $\lambda$ is the volume (Lebesgue measure) on $[0,1)^d$,
$I(F)=\int_{[0,1)^d}F(u)\,\dd u$, and for points $P=(x_1,\dots,x_n)$ of $[0,1)^d$
(repetitions allowed), $Q_P(F)=\frac1n\sum_{i=1}^nF(x_i)$ is the equal-weight rule. We write
$D=\{1,\dots,d\}$; for $u\subseteq D$, $x_u=(x_j)_{j\in u}$, and $(x_u{:}1)$ is the point whose
$j$-th coordinate is $x_j$ for $j\in u$ and $1$ for $j\notin u$.

\section{Discrepancy and the Koksma--Hlawka inequality}\label{sec:qmc-discrepancy}

Let $Z\sim\mathcal N(0,I_d)$ and let $f(Z)$ be a discounted payoff, with one coordinate per
date or per asset. If $U$ is uniform on $(0,1)^d$, the vector
$N^{-1}(U)=(N^{-1}(U_1),\dots,N^{-1}(U_d))$ has independent standard normal coordinates
(\cref{prop:rng-inversion}), so
\begin{equation}\label{eq:qmc-cube}
  \E[f(Z)]=\int_{(0,1)^d}F(u)\,\dd u,\qquad F=f\circ N^{-1}.
\end{equation}
Monte Carlo evaluates $Q_P(F)$ at independent uniform points; QMC chooses the points. The
dimension $d$ is the number of uniforms that one sample consumes, so it must be fixed:
samplers that consume a random number of uniforms, such as rejection, are excluded. The map
from uniforms to normals should also be monotone in each coordinate, as inversion is, so that
evenly spread uniforms give evenly spread normals.

\begin{definition}[Discrepancy]\label{def:qmc-discrepancy}
Let $P=(x_1,\dots,x_n)$ be points of $[0,1)^d$. For $a\in[0,1]^d$, let
$A_P(a)=\#\{i:x_i\in[0,a)\}$ count the points in the anchored box
$[0,a)=\prod_{j=1}^d[0,a_j)$, and let
\[
  \delta_P(a)=\frac{A_P(a)}{n}-\lambda\big([0,a)\big)=\frac{A_P(a)}{n}-\prod_{j=1}^da_j
\]
be the \emph{local discrepancy}. The \emph{star discrepancy} of $P$ is
$D^*_n(P)=\sup_{a\in[0,1]^d}|\delta_P(a)|$. For an infinite sequence, $D^*_n$ is the star
discrepancy of its first $n$ points. The sequence is \emph{uniformly distributed} if
$D^*_n\to0$, and it is a \emph{low-discrepancy sequence} if $D^*_n=O\big((\log n)^d/n\big)$.
\end{definition}

The local discrepancy is the integration error of $Q_P$ on the indicator of a box:
$\delta_P(a)=Q_P(\1_{[0,a)})-I(\1_{[0,a)})$. In one dimension,
$D^*_n(P)=\sup_x|\hat F_n(x)-x|$ is the Kolmogorov--Smirnov distance between the empirical
distribution function $\hat F_n$ of the points and the uniform law, the statistic of the test
used in \cref{ch:05-random-numbers}; for independent uniform points $\sqrt n\,D^*_n$ converges
in law to Kolmogorov's distribution (\citealp{knuth-taocp2}, Section~3.3.1.B), so random points
have discrepancy of order $n^{-1/2}$.
No point set does better than order $1/n$: in one dimension $D^*_n\ge\frac1{2n}$
(\exref{ex:qmc-one-dim}), and in dimension $d\ge2$ every set of $n$ points has
$D^*_n\ge c_d(\log n)^{(d-1)/2}/n$ for a constant $c_d>0$ \citep{roth-1954}. For infinite
sequences in $[0,1)$, $nD^*_n\ge c\log n$ for infinitely many $n$ \citep{schmidt-1972}.

\subsection*{Variation and the error bound}

The error bound needs a notion of variation. Call $f:[0,1]^d\to\R$ \emph{smooth} if all its
mixed partial derivatives $\partial^uf=\partial^{|u|}f/\prod_{j\in u}\partial x_j$,
$u\subseteq D$, exist and are continuous on the closed cube. For smooth $f$, the
\emph{variation in the sense of Hardy and Krause} is
\begin{equation}\label{eq:qmc-vhk}
  V_{HK}(f)=\sum_{\emptyset\ne u\subseteq D}\int_{[0,1]^{|u|}}\big|\partial^uf(a_u{:}1)\big|\,\dd a_u ,
\end{equation}
one term per face of the cube that contains the corner $(1,\dots,1)$. In one dimension it is
the total variation $\int_0^1|f'|$. The general definition, through sums of alternating
differences over grids, agrees with \eqref{eq:qmc-vhk} for smooth $f$ \citep{owen-2005}.

\begin{lemma}[Expansion from the corner]\label{lem:qmc-expansion}
If $f$ is smooth, then for every $x\in[0,1]^d$
\begin{equation}\label{eq:qmc-expansion}
  f(x)=\sum_{u\subseteq D}(-1)^{|u|}\int_{[0,1]^{|u|}}\1\{x_u<a_u\}\,\partial^uf(a_u{:}1)\,\dd a_u,
\end{equation}
where $\1\{x_u<a_u\}=\prod_{j\in u}\1\{x_j<a_j\}$ and the term $u=\emptyset$ is
$f(1,\dots,1)$. In particular $|f(x)-f(1,\dots,1)|\le V_{HK}(f)$.
\end{lemma}

\begin{proof}
For a function $g$ of one variable with a continuous derivative,
$g(x)=g(1)-\int_x^1g'(a)\,\dd a=g(1)-\int_0^1\1\{x<a\}\,g'(a)\,\dd a$. Apply this in $x_1$
to $f$, with the other coordinates fixed. Each of the two resulting terms is a function of
$x_2,\dots,x_d$ (integrated against $a_1$ in the second term) whose partial derivative in $x_2$
is continuous; apply the identity in $x_2$ to it, under the integral sign, and so on. After
the $d$-th step, each coordinate $j$ has either been set to $1$ or replaced by
$-\int_0^1\1\{x_j<a_j\}(\cdot)\,\dd a_j$ with one more derivative $\partial_j$, which is
\eqref{eq:qmc-expansion}. All integrands are continuous on compact sets, so the iterated
integrals are integrals over $[0,1]^{|u|}$ (Fubini). The bound follows by taking absolute
values in the terms $u\ne\emptyset$.
\end{proof}

\begin{theorem}[Koksma--Hlawka]\label{thm:qmc-koksma-hlawka}
Let $f:[0,1]^d\to\R$ have bounded variation $V_{HK}(f)$ in the sense of Hardy and Krause.
For all points $P=(x_1,\dots,x_n)$ of $[0,1)^d$,
\begin{equation}\label{eq:qmc-kh}
  \big|Q_P(f)-I(f)\big|\le V_{HK}(f)\,D^*_n(P).
\end{equation}
If $f$ is smooth, the error has the exact representation
\begin{equation}\label{eq:qmc-zaremba}
  Q_P(f)-I(f)=\sum_{\emptyset\ne u\subseteq D}(-1)^{|u|}\int_{[0,1]^{|u|}}
  \delta_P(a_u{:}1)\,\partial^uf(a_u{:}1)\,\dd a_u .
\end{equation}
\end{theorem}

The inequality is due to Koksma for $d=1$ and to Hlawka in general; see
\textcite{niederreiter-qmc}, Chapter~2. We prove it for smooth $f$,
which shows where the variation comes from.

\begin{proof}[Proof for smooth $f$]
Average \eqref{eq:qmc-expansion} over the points, and integrate it over $x$. For
$j\notin u$ the points satisfy $x_{ij}<1$, so $\frac1n\sum_i\1\{x_{i,u}<a_u\}=A_P(a_u{:}1)/n$,
while $\int\1\{x_u<a_u\}\,\dd x=\prod_{j\in u}a_j=\lambda([0,(a_u{:}1)))$. The terms $u=\emptyset$
cancel, and the difference of the other terms is \eqref{eq:qmc-zaremba} (the sums are
finite and the integrands bounded, so averaging commutes with the integrals). Each
$[0,(a_u{:}1))$ is an anchored box, so $|\delta_P(a_u{:}1)|\le D^*_n(P)$, and
\eqref{eq:qmc-kh} follows with $V_{HK}(f)$ given by \eqref{eq:qmc-vhk}.
\end{proof}

The bound separates the integrand from the points: the variation measures the integrand,
the discrepancy the points. In one dimension it cannot be improved: the indicator
$\1_{[0,a)}$ has variation $1$ and error $\delta_P(a)$, whose supremum over $a$ is
$D^*_n(P)$.

\begin{remark}[What the inequality says about option payoffs]\label{rem:qmc-payoffs}
Very little, for three reasons.
\begin{enumerate}[label=(\alph*)]
  \item A function of bounded Hardy--Krause variation is bounded: \cref{lem:qmc-expansion}
    gives $|f(x)-f(1,\dots,1)|\le V_{HK}(f)$ for smooth $f$, and the same inclusion--exclusion
    argument, with differences in place of integrals, gives it in general. The integrand
    $F=f\circ N^{-1}$ of \eqref{eq:qmc-cube} is unbounded as soon as the payoff is unbounded
    in $Z$ (the stock, a call, an Asian call), so $V_{HK}(F)=\infty$ and \eqref{eq:qmc-kh}
    is empty.
  \item Bounded payoffs are not safe either. For a hyperplane $a\cdot x=c$ that cuts the cube,
    a jump $\1\{a\cdot x>c\}$ has finite variation if only one $a_j$ is nonzero (a digital in
    one coordinate has variation $1$) and infinite variation as soon as two are; a kink
    $(a\cdot x-c)^+$ has finite variation if at most two $a_j$ are nonzero (its mixed differences
    in those two coordinates have the constant sign of their product and telescope) and infinite
    variation as soon as three are \citep{owen-2005}; the cases $\1\{x_1+x_2\le1\}$,
    $(x_1+x_2-1)^+$ and $(x_1+x_2+x_3-1)^+$ are \exref{ex:qmc-variation}. A digital on an average of two
    or more dates, or a put on an average of three or more, jumps or kinks across a curved surface
    of this kind after the map $N^{-1}$. A put on the average of two dates has infinite variation
    too, but because of $N^{-1}$ itself: with $u_j=N(z_j)$, the term $u=\{1,2\}$ of
    \eqref{eq:qmc-vhk} becomes $\int|\partial_{z_1}\partial_{z_2}f|\,\dd z_1\dd z_2$ with Lebesgue
    measure on $\R^2$, which diverges as $z_1\to-\infty$.
  \item Even for smooth integrands the bound is hard to use: $V_{HK}$ involves $2^d-1$
    integrals of mixed derivatives, computing $D^*_n$ exactly is infeasible beyond small $n$
    and $d$, and $(\log n)^d/n$ falls below $n^{-1/2}$ only when $n^{1/2}>(\log n)^d$, which
    for $d=10$ and $n\ge3$ means $n>10^{39}$.
\end{enumerate}
\end{remark}

What does follow without any variation is convergence.

\begin{proposition}[Convergence without a rate]\label{prop:qmc-riemann}
Let $x_1,x_2,\dots$ be uniformly distributed in $[0,1)^d$, and let $F:[0,1)^d\to\R$ be Riemann
integrable: for every $\varepsilon>0$ there are step functions $g\le F\le h$ (finite linear
combinations of indicators of boxes $[a,b)\subseteq[0,1)^d$) with $I(h-g)\le\varepsilon$. Then $Q_n(F)\to I(F)$,
where $Q_n$ uses the first $n$ points.
\end{proposition}

\begin{proof}
A box $B=[a,b)\subseteq[0,1)^d$ is a signed combination of $2^d$ anchored boxes, by
inclusion--exclusion: $\1_B=\sum_{v\subseteq D}(-1)^{|v|}\1_{[0,c^v)}$ with $c^v_j=a_j$ for
$j\in v$ and $c^v_j=b_j$ otherwise. Hence $|Q_n(\1_B)-\lambda(B)|\le2^dD^*_n\to0$, and
$Q_n(g)\to I(g)$ for every step function $g$. With $g\le F\le h$ as in the statement,
$\limsup_nQ_n(F)\le\lim_nQ_n(h)=I(h)\le I(F)+\varepsilon$ and likewise
$\liminf_nQ_n(F)\ge I(F)-\varepsilon$. Let $\varepsilon\downarrow0$.
\end{proof}

A bounded payoff whose discontinuities lie on a smooth surface, such as a digital on an
average after the map $N^{-1}$, is Riemann integrable: on a fine grid, the cells that meet the
closure of the surface have small total volume, and on the others the oscillation of $F$ is
small. So QMC converges for it, at a rate that no theorem gives.
For an unbounded $F\ge0$, uniform distribution alone does not even give convergence: replacing
the points of indices $n=2^k$ by points so close to a point where $F$ is unbounded that
$F(x_n)\ge n^2$ changes $D^*_n$ by
at most $(\log_2n+1)/n$, but makes $Q_n(F)\ge n$ along $n=2^k$. In one dimension a truncation argument
recovers a rate for monotone integrands (\exref{ex:qmc-unbounded}), and
\textcite{owen-2006} gives conditions on the growth of the mixed derivatives near the boundary
of the cube under which QMC keeps a rate close to $1/n$; they require mixed derivatives,
which kinked payoffs do not have. Randomised QMC (\cref{sec:qmc-rqmc}) needs only
$F\in L^2$.

\section{Low-discrepancy sequences}\label{sec:qmc-sequences}

\subsection{Van der Corput and Halton}

Let $b\ge2$ be an integer. Every integer $i\ge0$ has a unique expansion
$i=\sum_{k\ge0}a_k(i)\,b^k$ with digits $a_k(i)\in\{0,\dots,b-1\}$, finitely many of them
nonzero. The \emph{radical inverse} mirrors the digits at the radix point,
\[
  \phi_b(i)=\sum_{k\ge0}a_k(i)\,b^{-k-1}\in[0,1),
\]
and the \emph{van der Corput sequence} in base $b$ is $\phi_b(0)$, $\phi_b(1)$,
$\phi_b(2),\dots$; in base $2$ it starts $0$, $\frac12$, $\frac14$, $\frac34$, $\frac18$,
$\frac58$, $\frac38$, $\frac78$. The
\emph{Halton sequence} \citep{halton-1960} uses one base per coordinate,
$x_i=\big(\phi_{b_1}(i),\dots,\phi_{b_d}(i)\big)$, with pairwise coprime bases, usually the
first $d$ primes.

\begin{proposition}[Van der Corput and Halton]\label{prop:qmc-van-der-corput}
Let $b\ge2$.
\begin{enumerate}[label=(\roman*)]
  \item For every $k\ge0$, any $b^k$ consecutive points $\phi_b(i)$, $j\le i<j+b^k$, lie one
    in each interval $[l\,b^{-k},(l+1)\,b^{-k})$, $0\le l<b^k$. The first $b^k$ points are
    exactly the multiples of $b^{-k}$.
  \item For every $n\ge1$, the first $n$ points of the van der Corput sequence satisfy
    \[
      D^*_n\le\frac{(b-1)\big(\lfloor\log_bn\rfloor+1\big)}{n},
    \]
    and $D^*_n=1/n$ when $n$ is a power of $b$.
  \item Let $b_1,\dots,b_d$ be pairwise coprime and $k_1,\dots,k_d\ge0$. Any
    $B=\prod_jb_j^{k_j}$ consecutive points of the Halton sequence lie one in each box
    $\prod_j\big[l_jb_j^{-k_j},(l_j+1)b_j^{-k_j}\big)$.
\end{enumerate}
\end{proposition}

\begin{proof}
(i) The tail $\sum_{r\ge k}a_r(i)b^{-r-1}$ is smaller than $\sum_{r\ge k}(b-1)b^{-r-1}=b^{-k}$,
because the digits of an integer are eventually $0$. Hence the interval that contains
$\phi_b(i)$ has index $\lfloor b^k\phi_b(i)\rfloor=\sum_{r<k}a_r(i)\,b^{k-1-r}$: the first $k$
digits of $i$ in reverse order. It depends on $i$ only through $i\bmod b^k=\sum_{r<k}a_r(i)b^r$,
and digit reversal is a bijection of $\{0,\dots,b^k-1\}$. Consecutive integers $j,\dots,j+b^k-1$
take every residue modulo $b^k$ once. For $i<b^k$ the tail vanishes, so $\phi_b(i)=l\,b^{-k}$.

(ii) Let $K=\lfloor\log_bn\rfloor$ and $n=\sum_{k=0}^Kc_kb^k$ in base $b$. Split the indices
$0,\dots,n-1$ into consecutive blocks: $c_K$ blocks of length $b^K$, then $c_{K-1}$ of length
$b^{K-1}$, and so on down to $c_0$ blocks of length $1$; there are
$\sum_kc_k\le(b-1)(K+1)$ blocks. Fix $a\in[0,1]$ and a block of length $b^k$. By (i) it has
one point in each interval of length $b^{-k}$; $[0,a)$ contains the
$\lfloor ab^k\rfloor$ intervals to the left of $a$ and meets at most one more, so the block puts
$\lfloor ab^k\rfloor$ or $\lfloor ab^k\rfloor+1$ points in $[0,a)$, which differs from
$ab^k$ by less than $1$. Summing over blocks, $|A_P(a)-na|<(b-1)(K+1)$; divide by $n$. When
$n=b^k$ the points are $l/n$, $A_P(a)=\lceil an\rceil$, and $\delta_P(a)=(\lceil an\rceil-an)/n$
takes values in $[0,1/n)$ with supremum $1/n$.

(iii) By (i), the box that contains $x_i$ is determined by the residues $i\bmod b_j^{k_j}$,
$j=1,\dots,d$, through digit reversals, which are bijections. The moduli $b_j^{k_j}$ are pairwise
coprime, so by the Chinese remainder theorem $i\bmod B\mapsto(i\bmod b_j^{k_j})_{j\le d}$ is a
bijection. $B$ consecutive integers take every residue modulo $B$ once.
\end{proof}

Part (ii) is of the best possible order for a sequence in one dimension, by Schmidt's
theorem; sharper constants are in \textcite{niederreiter-qmc}, Chapter~3. In base $2$, the
largest value of $nD^*_n$ for $n\le4096$ is $4.78$, at $n=3413$, where (ii) allows $12$. Part
(iii), with the decomposition into blocks of (ii) in each coordinate, gives the bound
\[
  D^*_n<\frac dn+\frac1n\prod_{j=1}^d\Big(\frac{b_j-1}{2\log b_j}\log n+\frac{b_j+1}{2}\Big)
\]
for the Halton sequence \citep[Chapter~3]{niederreiter-qmc}: a low-discrepancy sequence,
but with a constant $\prod_j(b_j-1)/(2\log b_j)$ that grows superexponentially in $d$ when the
$b_j$ are the first $d$ primes. The reason is
visible in the coordinates with large bases. For $i<b$, $\phi_b(i)=i/b$: the coordinate
creeps up in $b$ equal steps. Two coordinates with nearby large bases creep up together, and
the first points of their projection lie on a few lines (\cref{fig:qmc-points}, right;
\exref{ex:qmc-halton-lines}). Randomly permuting the digits in each base, as the scramble of
\cref{sec:qmc-rqmc} does in base $2$, breaks the lines; shifting the points does not.

\begin{figure}[ht]
  \centering
  \includegraphics[width=\textwidth]{ch11-points}
  \caption{The first $256$ points of four sequences, projected on two coordinates.
    Pseudo-random points (numpy \texttt{default\_rng}, seed $0$) leave gaps and clusters. The
    Halton points in bases $2$ and $3$ and the Sobol' points in coordinates $1$ and $2$ fill
    the square evenly; the grey grid shows that each of the $16\times16$ squares of side
    $\frac1{16}$ holds exactly one Sobol' point. The Halton coordinates $15$ and $16$, with bases
    $47$ and $53$, lie on a few parallel lines.}
  \label{fig:qmc-points}
\end{figure}

\subsection{Nets and the Sobol' sequence}

Uniformity in all coordinates at once is easiest to organise in a single base. The
following notions are due to Sobol' and, in this generality, to Niederreiter
\citep[Chapter~4]{niederreiter-qmc}.

\begin{definition}[Nets and $(t,d)$-sequences]\label{def:qmc-net}
Let $b\ge2$ and $0\le t\le m$ be integers. An \emph{elementary interval} in base $b$ is a box
$E=\prod_{j=1}^d\big[l_jb^{-k_j},(l_j+1)b^{-k_j}\big)$ with integers $k_j\ge0$ and
$0\le l_j<b^{k_j}$; its volume is $b^{-\sum_jk_j}$. A \emph{$(t,m,d)$-net} in base $b$ is a set of
$n=b^m$ points of $[0,1)^d$ such that every elementary interval of volume $b^{t-m}$ contains
exactly $b^t$ points. A sequence $x_0,x_1,\dots$ is a \emph{$(t,d)$-sequence} in base $b$ if,
for every $m\ge t$ and $q\ge0$, the points $x_i$ with $qb^m\le i<(q+1)b^m$ form a
$(t,m,d)$-net.
\end{definition}

A net with $t=0$ is balanced in every elementary interval of volume $1/n$: every square of
side $\frac1{16}$ in \cref{fig:qmc-points}, but also every rectangle of $\frac12\times\frac1{128}$
or $\frac1{256}\times1$. By \cref{prop:qmc-van-der-corput}(i), the van der Corput sequence in
base $b$ is a $(0,1)$-sequence in base $b$.

\begin{definition}[Sobol' sequence]\label{def:qmc-sobol}
For each coordinate $j\ge2$, take a primitive polynomial over the field $\mathbb F_2$ of
degree $s=s_j$, $p_j(z)=z^s+c_1z^{s-1}+\dots+c_{s-1}z+1$, a different one for each coordinate,
and odd integers $m_1,\dots,m_s$
with $m_k<2^k$. Extend them by
\[
  m_k=2c_1m_{k-1}\oplus2^2c_2m_{k-2}\oplus\dots\oplus2^{s-1}c_{s-1}m_{k-s+1}\oplus2^sm_{k-s}
  \oplus m_{k-s},\qquad k>s,
\]
where $\oplus$ is the bitwise exclusive or of binary expansions, and define the
\emph{direction numbers} $v_{k,j}=m_k/2^k$. For coordinate $1$, $s_1=1$ and $m_k=1$ for all $k$.
The $i$-th point of the \emph{Sobol' sequence}, for $i=\sum_{k\ge1}i_k2^{k-1}$ with bits
$i_k\in\{0,1\}$, has coordinates
\[
  x_{i,j}=i_1v_{1,j}\oplus i_2v_{2,j}\oplus i_3v_{3,j}\oplus\cdots,\qquad j=1,\dots,d.
\]
\end{definition}

The recursion keeps every $m_k$ odd and smaller than $2^k$: the term $m_{k-s}$ is odd, all
the others are even, and each is smaller than $2^k$. Coordinate $1$ is the van der Corput
sequence in base $2$, since $x_{i,1}=\sum_ki_k2^{-k}=\phi_2(i)$. The construction is due to
\textcite{sobol-1967}; its quality depends on the choice of the initial numbers $m_k$. We use
those of \textcite{joe-kuo-2008}, chosen to improve the two-dimensional projections, for the
first $128$ coordinates (see also \textcite{glasserman-mc}, Section~5.2.3).

\begin{proposition}[Sobol' points form digital nets]\label{prop:qmc-digital-net}
Fix $m\ge1$, and for each coordinate $j$ let $C_j$ be the $m\times m$ matrix over
$\mathbb F_2$ whose $k$-th column holds the first $m$ binary digits of $v_{k,j}$.
\begin{enumerate}[label=(\roman*)]
  \item For $0\le i<2^m$, $x_{i,j}$ has at most $m$ binary digits, and they are
    $C_j(i_1,\dots,i_m)^\top$ (arithmetic modulo $2$). Each $C_j$ is upper triangular with
    ones on the diagonal. Consequently, in every block $q2^m\le i<(q+1)2^m$, each coordinate
    has one point in each interval $[l2^{-m},(l+1)2^{-m})$.
  \item Every block $q2^m\le i<(q+1)2^m$ is a $(t,m,d)$-net in base $2$ if and only if, for every
    $(k_1,\dots,k_d)$ with $k_j\ge0$ and $\sum_jk_j=m-t$, the first $k_1$ rows of $C_1$, the first
    $k_2$ rows of $C_2$, \dots, and the first $k_d$ rows of $C_d$ are linearly independent over
    $\mathbb F_2$.
  \item If the polynomials $p_j$ are distinct, the Sobol' sequence is a $(t,d)$-sequence in base
    $2$ with $t\le\sum_{j=1}^d(s_j-1)$.
\end{enumerate}
\end{proposition}

\begin{proof}
(i) The bitwise exclusive or adds binary digits modulo $2$ without carries, so digit $r$ of
$x_{i,j}$ is $\sum_ki_k\,(\text{digit $r$ of }v_{k,j})$ modulo $2$. Since $m_k$ is odd and
$m_k<2^k$, the digits of $v_{k,j}=m_k/2^k$ beyond the $k$-th are $0$ and the $k$-th is $1$:
column $k$ of $C_j$ vanishes below row $k$ and has a $1$ in row $k$. For $i<2^m$ only
$i_1,\dots,i_m$ can be nonzero, which gives the first claim. In the block of index $q$, write
$i=i'+q2^m$ with $i'<2^m$; the bits of $q2^m$ contribute the same vector $e_j$ to the first $m$
digits of every point, so those digits are $C_j(i'_1,\dots,i'_m)^\top+e_j$. As $i'$ runs over
$\{0,\dots,2^m-1\}$, this runs over all $2^m$ digit strings, because $C_j$ is invertible.

(ii) A point lies in the elementary interval with parameters $(k_j,l_j)$ if and only if, for
each $j$, the first $k_j$ digits of its $j$-th coordinate are the $k_j$ digits of $l_j$. By (i)
this is a system of $\sum_jk_j=m-t$ linear equations over $\mathbb F_2$ in the $m$ unknown bits of
$i'$, whose matrix $M$ stacks the first $k_j$ rows of each $C_j$, and whose right-hand side is the
vector of digits of the $l_j$ minus the shift of (i); as the $l_j$ vary, it runs over all of
$\mathbb F_2^{m-t}$. If the rows of $M$ are independent, $M$ is onto and every system has exactly
$2^{m-(m-t)}=2^t$ solutions. If they are dependent, $M$ is not onto, so some right-hand side has no
solution: some elementary interval of volume $2^{t-m}$ is empty, and the block is not a
$(t,m,d)$-net.

(iii) is due to \textcite{sobol-1967}; see \textcite{dick-pillichshammer}, Chapter~8.
\end{proof}

By (ii), the exact $t$-value of the first $2^m$ points is the smallest $t$ that passes the rank
test, which the code computes by Gaussian elimination over $\mathbb F_2$ and checks by counting
points in every elementary interval. For the first $2^{10}$ points, the
first $2$, $3$, $5$ and $8$ coordinates have $t=0$, $1$, $3$ and $5$, against the bounds $0$, $1$,
$5$ and $15$ of (iii). Nets have low discrepancy: for a $(t,m,d)$-net in base $b$,
$D^*_n\le C_{d,b}\,b^t(\log n)^{d-1}/n+O\big(b^t(\log n)^{d-2}/n\big)$, and for the first $n$ points
of a $(t,d)$-sequence, for every $n\ge2$,
$D^*_n\le C'_{d,b}\,b^t(\log n)^d/n+O\big(b^t(\log n)^{d-1}/n\big)$
\citep[Chapter~4; Theorem~4.17 for sequences]{niederreiter-qmc}. In two dimensions the star discrepancy can be
computed exactly, and at $n=2048$ it is $0.020$ for pseudo-random points, $0.0037$ for Halton and
$0.0024$ for Sobol' (\cref{fig:qmc-convergence}, left).

\begin{remark}[The origin, and $n$ not a power of two]\label{rem:qmc-origin}
Both sequences start at $x_0=0$, which $N^{-1}$ sends to $-\infty$. A deterministic run with
the normal map must drop the first point, which spoils the net property of the first $2^m$
points, or move all the points: by (i) the first $2^m$ Sobol' points have coordinates that are
multiples of $2^{-m}$, and adding $2^{-m-1}$ to every coordinate keeps each point in all the
elementary intervals of side at least $2^{-m}$ that contained it, so the result is still a
$(t,m,d)$-net, with no coordinate at $0$. Randomisation settles the issue. Second, the balance of
\cref{def:qmc-net} concerns blocks of $2^m$ points with $m\ge t$. For other $n$, the first $n$
points split into aligned blocks of sizes $2^k$ (as in the proof of
\cref{prop:qmc-van-der-corput}(ii)), which are $(t,k,d)$-nets only for $k\ge t$; the blocks with
$k<t$ carry no such guarantee. The discrepancy bound for sequences still holds, but the faster
rates of \cref{sec:qmc-rqmc} are lost (\cref{fig:qmc-convergence}, right).
\end{remark}

\begin{incode}
\coderef{code/src/mfe230d/qmc.py}: \texttt{van\_der\_corput(n, b)}, \texttt{halton(n, d)} and
\texttt{sobol(n, d)} (natural order, direction numbers of \citeauthor{joe-kuo-2008} for up to
$128$ coordinates; \texttt{scipy.stats.qmc.Sobol} returns the same points in Gray-code order,
which the tests check); \texttt{star\_discrepancy} (exact, practical for $d\le2$);
\texttt{t\_value(generator\_matrices(d, m))} and \texttt{is\_net}.
\end{incode}

\section{Randomised quasi-Monte Carlo}\label{sec:qmc-rqmc}

A deterministic rule has no standard error, and its error bound is rarely usable
(\cref{rem:qmc-payoffs}). Randomisation makes every point uniform while keeping the
structure of the set. Three randomisations of points $x_1,\dots,x_n$ are in use; each draws its
randomness once and applies it to all the points.
\begin{enumerate}[label=(\alph*)]
  \item \emph{Random shift} \citep{cranley-patterson}: $y_i=x_i+U\bmod1$ (coordinatewise
    fractional part), with $U$ uniform on $[0,1)^d$.
  \item \emph{Random digital shift} in base $2$: $y_i=x_i\oplus U$, digit by digit.
  \item \emph{Nested uniform scramble} \citep{owen-1995}, in base $2$: in each coordinate,
    independently, digit $k$ of $y$ is digit $k$ of $x$ flipped by a random bit
    $e(x^{(1)},\dots,x^{(k-1)})$ attached to the first $k-1$ digits of $x$; all the bits $e$ are
    independent fair coin tosses. Geometrically: swap the two halves of $[0,1)$ with probability
    $\frac12$, then independently swap the two halves of each half, and so on.
\end{enumerate}

\begin{proposition}[Randomised QMC]\label{prop:qmc-rqmc}
Let $y_1,\dots,y_n$ be obtained from $x_1,\dots,x_n\in[0,1)^d$ by (a), (b) or (c), let $F\in
L^2([0,1)^d)$ with $\sigma^2=\Var F(U)$ for $U$ uniform, and let $\hat\mu=\frac1n\sum_iF(y_i)$.
\begin{enumerate}[label=(\roman*)]
  \item Each $y_i$ is uniformly distributed on $[0,1)^d$. Hence $\E\hat\mu=I(F)$, and
    $\Var\hat\mu\le\sigma^2$.
  \item The bound $\Var\hat\mu\le\sigma^2$, $n$ times the Monte Carlo variance, is attained for
    some point sets and integrands under (a) and (b), and under (c) when all the points coincide.
  \item Under (b) and (c), if $x_1,\dots,x_n$ is a $(t,m,d)$-net in base $2$, then so is
    $y_1,\dots,y_n$, with probability one.
  \item Let $\hat\mu^{(1)},\dots,\hat\mu^{(R)}$, $R\ge2$, be computed from independent
    randomisations, with mean $\bar\mu_R$ and sample variance $s_R^2$. Then $\bar\mu_R$ is
    unbiased, $\Var\bar\mu_R=\Var\hat\mu/R$, and $\E[s_R^2/R]=\Var\bar\mu_R$. If the replicates
    are normally distributed, $\bar\mu_R\pm t_{R-1,1-\alpha/2}\,s_R/\sqrt R$ covers $I(F)$ with
    probability exactly $1-\alpha$; without normality, its coverage tends to $1-\alpha$ as
    $R\to\infty$.
  \item Let $x_1,\dots,x_n$ be a $(t,m,d)$-net in base $2$, $n=2^m$, randomised by (c). Then
    $\Var\hat\mu\le2^t3^d\sigma^2/n$. For the first $2^m$ points of a $(t,d)$-sequence,
    $\Var\hat\mu=o(1/n)$ as $m\to\infty$, and if moreover $F$ is smooth,
    $\Var\hat\mu=O\big(n^{-3}(\log n)^{d-1}\big)$.
\end{enumerate}
\end{proposition}

\begin{proof}
(i) Under (a), $u\mapsto x+u\bmod1$ is a translation of the torus and preserves volume. Under
(b) and (c), recall that $U$ is uniform on $[0,1)$ if and only if its binary digits are
independent fair bits (outside the null set of dyadic rationals). Under (b), each digit of $y$
is a fixed digit of $x$ flipped by an independent fair bit; under (c), digit $k$ of $y$ is
digit $k$ of $x$ flipped by the bit attached to one node at depth $k-1$, and different depths
use different, independent bits. In both cases the digits of each coordinate of $y$ are
independent fair bits, and different coordinates are independent. So
$\E F(y_i)=I(F)$, and by the Cauchy--Schwarz inequality
$\Var\hat\mu=n^{-2}\sum_{i,k}\Cov(F(y_i),F(y_k))\le n^{-2}\sum_{i,k}\sigma^2=\sigma^2$.

(ii) For (a), take $d=1$, $x_i=(i-1)/n$ and $F(y)=\cos(2\pi ny)$: every
$F(y_i)=\cos(2\pi nU)$, so $\hat\mu=\cos(2\pi nU)$ has variance $\frac12=\sigma^2$. For (b), take
$n=2^m$ and $F(y)=(-1)^{y^{(m+1)}}$, the sign given by digit $m+1$: that digit is $0$ for every
$x_i=(i-1)/n$, so it equals $U^{(m+1)}$ for every $y_i$, and $\hat\mu=\pm1$ with variance
$1=\sigma^2$.
If all points coincide, so do the $y_i$, and $\hat\mu=F(y_1)$.

(iii) Under (b) and (c), the first $k$ digits of each coordinate of $y$ are a bijective function
of the first $k$ digits of the same coordinate of $x$ (for (c), the composition of the flips
along the path from the root). So each elementary interval is mapped onto an elementary
interval of the same shape, and the counts of \cref{def:qmc-net} are preserved. The digit
expansions involved are proper with probability one.

(iv) The replicates are independent copies of $\hat\mu$, which is in $L^2$ by (i). Unbiasedness
of $\bar\mu_R$ and of the sample variance for independent identically distributed variables give
the first three claims; the normal case is the classical Student interval, and the limit as
$R\to\infty$ is \cref{prop:mc-batch} with the replicates as batches.

(v) The rate $o(1/n)$ for every $F\in L^2$ is due to \textcite{owen-scrambled-nets}, as is a
bound by a small multiple of $\sigma^2/n$ for $t=0$; the bound $2^t3^d\sigma^2/n$ is
\textcite{owen-1998}, Lemma~4. The rate for smooth $F$ is due to \textcite{owen-1997b};
\textcite{owen-2008}, Section~5.1 and Theorem~3, repairs an error in its proof and weakens its
hypothesis to the smoothness assumed here.
\end{proof}

Three comments. First, $\hat\mu$ is a sum of dependent terms, and the replicates of (iv), not the
spread of the $F(y_i)$, measure its error: this is the batch-means principle of
\cref{prop:mc-batch}. For a fixed number of replicates, the Student interval rests on
approximate normality of each replicate, which \textcite{loh-2003} proved for scrambled
$(0,m,d)$-nets and smooth integrands. Second, the budget: with $N=nR$ evaluations and
$\Var\hat\mu\approx cn^{-\alpha}$, $\Var\bar\mu_R\approx cn^{1-\alpha}/N$, so for $\alpha>1$ one
should take $n$ large and $R$ as small as a stable error estimate allows, say $10$ to $30$.
Third, in one dimension the nested scramble of the first $2^m$ van der Corput points is a
stratified sample with $2^m$ strata (\cref{ch:10-stratification}): one point in each interval
$[l2^{-m},(l+1)2^{-m})$, uniform there and independent of the others
(\exref{ex:qmc-scrambled-vdc}). In $d$ dimensions, by \cref{prop:qmc-digital-net}(i), each
one-dimensional projection of the scrambled first $2^m$ Sobol' points is stratified in this way, as
in Latin hypercube sampling (\cref{prop:st-lhs}), and the net also balances the elementary
intervals in several coordinates at once.

\begin{figure}[ht]
  \centering
  \includegraphics[width=\textwidth]{ch11-convergence}
  \caption{Left: exact star discrepancy of the first $n$ points in two dimensions, for
    pseudo-random points (seed $1$), Halton (bases $2$, $3$) and Sobol' (coordinates $1$, $2$),
    with reference slopes $n^{-1/2}$ and $\log n/n$. Centre: the product of sines
    $\prod_{j=1}^5\frac\pi2\sin(\pi x_j)$ of \cref{ch:06-monte-carlo}, whose integral is $1$:
    root mean squared error $\sigma/\sqrt n$ of Monte Carlo, absolute errors of the Halton and
    Sobol' rules, and the root mean squared error of nested-scrambled Sobol' points over $128$
    independent scrambles (seed $7$). Right: $n$ times the root mean squared error of scrambled
    Sobol' points for the same integrand in $d=2$, over $64$ scrambles (seed $8$), for every $n$
    from $16$ to $4096$; the dots mark $n=2^m$.}
  \label{fig:qmc-convergence}
\end{figure}

\Cref{fig:qmc-convergence} shows the three regimes. For the smooth integrand in $d=5$, at
$n=2^{16}$, Monte Carlo has root mean squared error $5.3\times10^{-3}$, the scrambled Sobol'
points $8.8\times10^{-5}$, and the deterministic Halton and Sobol' rules have errors
$2.8\times10^{-5}$ and $1.9\times10^{-4}$ (our implementation: numpy \texttt{default\_rng}, seeds as
in the caption). The deterministic errors jump around with $n$; the randomised one decreases
steadily, with a fitted slope of $-0.99$ in $\log$--$\log$ scale between $n=2^8$ and $2^{16}$
(from $-1.03$ to $-0.95$ over the four groups of $32$ scrambles; the replicate errors are
heavy-tailed, so fewer scrambles give noticeably different values):
the rate $n^{-3/2}$ of (v) is asymptotic, and its logarithmic factor $(\log n)^{(d-1)/2}$ still
grows fast at these $n$ when $d=5$. In $d=2$ the asymptotic regime is visible:
$n\times\text{RMSE}$ falls from $0.36$ at $n=2^8$ to $0.10$ at $n=2^{12}$ along powers of two,
while between them it is of order $1$ (median $1.2$ for $2^{11}<n<2^{12}$), the order $1/n$
that the smallest blocks of \cref{rem:qmc-origin} impose.

\demo{11-quasi-monte-carlo.ipynb}{scramble the Sobol' points and watch the net survive; then
compare Monte Carlo, QMC and scrambled QMC on smooth, kinked and discontinuous integrands, in the
dimension of your choice, with replicate-based error bars.}

\section{Effective dimension and path constructions}\label{sec:qmc-effective}

A monthly Asian option is an integral in $12$ dimensions, a weekly one in $52$. The bounds of
\cref{sec:qmc-sequences} are useless there, and yet QMC prices such options remarkably well.
The explanation of \textcite{caflisch-morokoff-owen} is that the integrands are
\emph{effectively} low-dimensional, provided the normals are mapped to the path in the right
order. It rests on the analysis-of-variance (ANOVA) decomposition. Every $F\in
L^2([0,1)^d)$ can be written uniquely as $F=\sum_{u\subseteq D}F_u$, where $F_u$ depends only
on $x_u$, $F_\emptyset=I(F)$, and $\int_0^1F_u\,\dd x_j=0$ for $j\in u$; the terms are
orthogonal, so $\sigma^2=\sum_{u\ne\emptyset}\sigma_u^2$ with $\sigma_u^2=I(F_u^2)$, and
$\E[F(U)\mid U_v]=\sum_{u\subseteq v}F_u(U_u)$ for every $v\subseteq D$
\citep{sobol-2001}. In particular
\begin{equation}\label{eq:qmc-closed-share}
  S_v:=\frac{\Var\E[F(U)\mid U_v]}{\sigma^2}=\frac1{\sigma^2}\sum_{\emptyset\ne u\subseteq v}\sigma_u^2
\end{equation}
is the share of the variance carried by the coordinates in $v$ and their interactions.

\begin{lemma}[Pick and freeze]\label{lem:qmc-pick-freeze}
Let $U,U'$ be independent and uniform on $[0,1)^d$, $v\subseteq D$, and let
$U^{(v)}=(U_v,U'_{D\setminus v})$ keep the coordinates of $U$ in $v$ and take the others from
$U'$. For $F\in L^2$, $\Var\E[F(U)\mid U_v]=\Cov\big(F(U),F(U^{(v)})\big)$.
\end{lemma}

\begin{proof}
Given $U_v$, the variables $F(U)$ and $F(U^{(v)})$ are independent with the same conditional
law, because they are the same function of $U_v$ and of the independent vectors
$U_{D\setminus v}$ and $U'_{D\setminus v}$. With $h(U_v)=\E[F(U)\mid U_v]$, the tower property
gives $\E[F(U)F(U^{(v)})]=\E[h(U_v)^2]$; subtract $I(F)^2=(\E h(U_v))^2$.
\end{proof}

The lemma turns each share into a covariance that plain Monte Carlo estimates
\citep{sobol-2001}. In normal coordinates it reads the same with $Z,Z'$ in place of $U,U'$.

\begin{remark}[Effective dimension]\label{rem:qmc-effective-dimension}
Following \textcite{caflisch-morokoff-owen}, with a proportion $p$ such as $0.99$, $F$ has
\emph{truncation dimension} $d_T$, the smallest $k$ with $S_{\{1,\dots,k\}}\ge p$, and
\emph{superposition dimension} $d_S$, the smallest $s$ with
$\sum_{0<|u|\le s}\sigma_u^2\ge p\sigma^2$. The first says how many leading coordinates matter,
the second how many coordinates interact; $\sum_jS_{\{j\}}$, the share of the additive part, is
the case $s=1$. Both depend on how the uniforms are turned into the path, although the price does
not: the increments, Brownian-bridge and principal-component constructions of
\cref{ch:07-path-simulation} give three integrands $F$ on $[0,1)^d$ with the same integral. The
bridge and principal components put most of the variance of the path into the first normals
(\cref{prop:ps-bridge,prop:ps-pca}), and QMC points are best in their first coordinates: the
leading Sobol' coordinates use the polynomials of lowest degree, hence the smallest bounds on $t$
(\cref{prop:qmc-digital-net}(iii)), and the leading Halton coordinates the smallest bases. For scrambled nets, part of this argument can be
made exact. The variance of (c) splits as $\Var\hat\mu=n^{-1}\sum_{u\ne\emptyset}\sum_r
\Gamma_{u,r}\sigma^2_{u,r}$, where the $\sigma^2_{u,r}\ge0$ split each $\sigma_u^2$ further by
resolution $r$ and the gains $\Gamma_{u,r}$ depend only on the points \citep{owen-scrambled-nets};
for a $(t,m,d)$-net in base $2$, $0\le\Gamma_{u,r}\le2^t3^d$ \citep{owen-1998}. The terms with $u\subseteq\{1,\dots,k\}$ make up the scrambled-net variance for
the $k$-dimensional integrand $\E[F\mid U_1,\dots,U_k]$, which is $o(1/n)$ by
\cref{prop:qmc-rqmc}(v); the others add at most $2^t3^d(1-S_{\{1,\dots,k\}})\sigma^2/n$. This much
is a theorem. That a small $1-S_{\{1,\dots,k\}}$ makes QMC accurate at a given $n$ is a
heuristic: nothing says how large $n$ must be before the $k$-dimensional part is integrated
accurately, and the constant $2^t3^d$ is astronomically large for $d=64$. It is the measurements
below that support it.
\end{remark}

For the \emph{geometric} average the shares are exact. On the dates $t_1<\dots<t_m$, with
$W=AZ$ for the construction matrix $A$,
\[
  \log G=\log S+\nu\bar t+\frac\sigma m\sum_{i=1}^mW_{t_i}
  =\log S+\nu\bar t+\frac\sigma m\,\big|A^\top\mathbf 1\big|\,(w\cdot Z),\qquad
  w=\frac{A^\top\mathbf 1}{|A^\top\mathbf 1|},
\]
with $\bar t$ the mean date. So the geometric Asian call is $g(w\cdot Z)$ for a fixed
function $g$ and a unit vector $w$ that depends on the construction. For such a \emph{ridge}
payoff, $(w\cdot Z,w\cdot Z^{(v)})$ is a standard bivariate normal pair with correlation
$\rho_v^2=\sum_{j\in v}w_j^2$. Write $\kappa_g(c)=\Corr\big(g(X),g(X')\big)$ for standard normal
$X,X'$ with $\Corr(X,X')=c$; then \cref{lem:qmc-pick-freeze} gives
\begin{equation}\label{eq:qmc-ridge}
  S_v=\kappa_g\big(\rho_v^2\big),\qquad \rho_v^2=\sum_{j\in v}w_j^2 .
\end{equation}
For a linear payoff $\kappa_g(c)=c$; in general $\kappa_g(c)\le c$ (\exref{ex:qmc-ridge}). For
the call, $\E[g(X')\mid X]$ is a Black--Scholes-type formula and $\kappa_g$ is a one-dimensional
integral. \Cref{tab:qmc-shares} gives the result for $m=64$ equally spaced dates over one year
in the test case.

\begin{table}[ht]
  \centering
  \small
  \begin{tabular}{@{}lcccccccc@{}}
    \toprule
    construction & $w_1^2$ & $S_{\{1\}}$ & $S_{\{1,2\}}$ & $S_{\{1,\dots,4\}}$ & $S_{\{1,\dots,8\}}$
      & $d_T$ & additive share \\
    \midrule
    increments & $0.046$ & $0.032$ & $0.063$ & $0.127$ & $0.253$ & $53$ & $0.684$\\
    Brownian bridge & $0.756$ & $0.685$ & $0.914$ & $0.978$ & $0.994$ & $7$ & $0.860$\\
    principal components & $0.986$ & $0.979$ & $0.997$ & $1.000$ & $1.000$ & $2$ & $0.989$\\
    \bottomrule
  \end{tabular}
  \caption{Variance shares of the geometric Asian call with $m=64$ dates (test case), by
    \eqref{eq:qmc-ridge} and quadrature: the linear share $w_1^2$ of the first normal, the
    shares $S_v$ of the first $k$ normals, the truncation dimension $d_T$ for $p=0.99$, and the
    share $\sum_jS_{\{j\}}$ of the additive part.}
  \label{tab:qmc-shares}
\end{table}

The first normal of the bridge sets $W_T$, which carries about $\frac34$ of the variance of
the average; the first principal component carries $98.6\%$ of it. The call loses a little
of each share to nonlinear interactions, much more under the increments, where the average
mixes all $64$ normals with comparable weights and only $68\%$ of the variance is additive.
The arithmetic average is not a ridge payoff, but its shares, estimated by
\cref{lem:qmc-pick-freeze}, are close to those of the geometric call (\cref{fig:qmc-asian},
right).

\begin{figure}[ht]
  \centering
  \includegraphics[width=\textwidth]{ch11-asian}
  \caption{The arithmetic Asian call with $m=64$ dates, test case. Left: standard deviation of
    the estimate with $n$ nested-scrambled Sobol' points, over $64$ independent scrambles
    (seed $11$), under the three path constructions, against the Monte Carlo standard error
    $\sigma/\sqrt n$ ($\sigma=7.12$, from the second moment of the payoff over all $2^{20}$
    scrambled points). Right: share of the variance carried by the first $k$
    normals: exact for the geometric Asian call (lines, \eqref{eq:qmc-ridge}), and estimated by
    pick-freeze for the arithmetic one (circles, $2^{15}$ draws, seed $5$).}
  \label{fig:qmc-asian}
\end{figure}

\begin{example}[Asian options by randomised QMC]\label{exa:qmc-asian}
With $m=64$ dates, $n=2^{14}$ scrambled Sobol' points and $64$ scrambles (our implementation,
seed $11$), the variance of the estimate is smaller than the Monte Carlo variance $\sigma^2/n$
with the same $n$ by a factor of $33$ with increments, $2.2\times10^3$ with the bridge and
$3.0\times10^4$ with principal components (\cref{fig:qmc-asian}, left). The replicate errors are
heavy-tailed, and the factors are noisy: over the four groups of $16$ scrambles they range from
$24$ to $50$, from $1.4\times10^3$ to $3.8\times10^3$, and from $1.5\times10^4$ to $6.2\times10^4$.
From $n=2^8$ to $2^{14}$ the factor of the bridge grows from $1.5\times10^2$ and that of principal
components from $6.0\times10^2$, while that of the increments levels off around $30$. The ranking
is that of \cref{tab:qmc-shares}. Principal components cost $O(m^2)$
operations per path against $O(m)$ for the other two, which does not change it here. For the course's monthly example ($m=12$), $n=2^{14}$
points and $16$ scrambles (seed $2026$) give $4.56133\pm0.00011$ for the arithmetic Asian call
with principal components and $4.56111\pm0.00038$ with the bridge, and $4.40291\pm0.00011$ for the
geometric one, whose exact value is $4.40286$. The course's estimate $4.56058$ came from $10^8$
plain paths, with a standard error of $7.5\times10^{-4}$, seven times that of principal components
with $2.6\times10^5$ paths; the two estimates differ by one of its standard errors.
\end{example}

The constructions are not uniformly ranked. A payoff that depends only on the first date,
$W_{t_1}$, is one-dimensional under the increments and spread over several normals by the
bridge (\exref{ex:qmc-first-date}); \textcite{papageorgiou-bridge} constructs integrands for
which the bridge does worse than the increments. The effective dimension belongs to the pair
(payoff, construction), and it can be measured before pricing with \cref{lem:qmc-pick-freeze}.
The use of the bridge for QMC is due to \textcite{moskowitz-caflisch}, and that of principal
components to \textcite{acworth-broadie-glasserman}.

\begin{incode}
\coderef{code/src/mfe230d/qmc.py}: \texttt{random\_shift}, \texttt{digital\_shift} and
\texttt{owen\_scramble} randomise points; \texttt{rqmc\_estimate(f, n, d, R)} returns the
replicate mean, its standard error and a Student interval; \texttt{to\_normals} applies
$N^{-1}$ (the accurate quantile of \cref{ch:05-random-numbers}); \texttt{asian\_call\_payoffs}
builds the paths with \texttt{paths.gbm\_paths}; \texttt{lognormal\_call\_share},
\texttt{ridge\_shares} and \texttt{pick\_freeze\_shares} compute the variance shares.
\end{incode}

\demo{11-quasi-monte-carlo.ipynb}{choose the number of dates and the construction, and compare the
variance shares of the first normals with the accuracy of randomised QMC for the Asian call.}

\paragraph{Course examples reproduced.} The course announced low-discrepancy sequences without
covering them, so it has no worked examples on this topic. We reproduce the numbers of its Asian
example that the chapter touches: the exact geometric Asian call $4.40286$, and the bare-bones
estimate $4.56058$ of the arithmetic one, which agrees with our randomised QMC value within its
standard error.

\furtherreading
\textcite{glasserman-mc}, Chapter~5, covers QMC for finance: discrepancy and the Koksma--Hlawka
bound in Section~5.1, the Halton, Faure and Sobol' constructions in Section~5.2, lattice rules,
which we have skipped, in Section~5.3, randomised QMC in Section~5.4, and numerical examples with
the bridge and principal components in Section~5.5. \textcite{niederreiter-qmc}, Chapters~2--4,
is the classical treatment of discrepancy, error bounds and $(t,m,s)$-nets, and
\textcite{dick-pillichshammer} the modern one of digital nets, including their randomisation in
Chapter~13. \textcite{owen-2005} explains the Hardy--Krause variation and why kinks and jumps
defeat it.
